(a) First suppose $X$ is connected. Since a regular local ring is a domain, the irreducible components of a regular noetherian scheme are disjoint and open. Thus $X$ is integral. Write $K=K(X)$ and $\pi:\mathbb P(\mathcal E)\to X$. Both schemes are locally factorial.

An invertible sheaf on $\mathbb P(\mathcal E)$ restricts to $\mathcal O(n)$ on $\mathbb P(\mathcal E)_K$ for a unique $n\in\mathbb Z$. After tensoring with $\mathcal O(-n)$ and subtracting the divisor of a rational function, its divisor has only vertical components. Every vertical prime divisor is $\pi^{-1}(D)$ for a prime divisor $D$ of $X$: in a polynomial chart its height-one prime contracts to a nonzero prime of height at most one and equals the extension of that prime. Since $X$ is locally factorial, these divisors are Cartier. Hence every invertible sheaf has the form $\pi^*\mathcal L\otimes\mathcal O(n)$. Uniqueness follows by restriction to the generic fibre and from $\pi_*\pi^*\mathcal L\cong\mathcal L$, using (7.11) and \exref{II.5.1}(d). Therefore $\operatorname{Pic}\mathbb P(\mathcal E)\cong\operatorname{Pic}X\oplus\mathbb Z$.

For arbitrary $X$, this argument applies to each component and gives the precise formula $\operatorname{Pic}\mathbb P(\mathcal E)\cong\operatorname{Pic}X\oplus\mathbb Z^{\pi_0(X)}$. Thus the stated formula requires connectedness. The argument also applies when $X$ is locally factorial.

(b) If $h:\mathbb P(\mathcal E)\to\mathbb P(\mathcal E')$ is an $X$-isomorphism, its restriction to every fibre pulls $\mathcal O(1)$ back to $\mathcal O(1)$ by (7.1.1). Part (a) therefore gives $h^*\mathcal O_{\mathbb P(\mathcal E')}(1)\cong\mathcal O_{\mathbb P(\mathcal E)}(1)\otimes\pi^*\mathcal L$ for some invertible $\mathcal L$ on $X$. Applying $\pi_*$ and (7.11) yields $\mathcal E'\cong\mathcal E\otimes\mathcal L$. Conversely, such an isomorphism gives $\mathbb P(\mathcal E')\cong\mathbb P(\mathcal E)$ by (7.9).
